% Generated draft: requires mathematical and source review.
\documentclass[11pt]{article}
\usepackage[margin=25mm]{geometry}
\usepackage{fontspec}
\setmainfont{Latin Modern Roman}
\usepackage{amsmath,amssymb,mathtools,unicode-math}
\setmathfont{Latin Modern Math}
\usepackage{enumitem,needspace,xurl,hyperref}
\setlength{\parindent}{0pt}
\setlength{\parskip}{5pt}
\setlength{\emergencystretch}{3em}
\newcommand{\N}{\mathbb{N}}
\newcommand{\Z}{\mathbb{Z}}
\newcommand{\Q}{\mathbb{Q}}
\newcommand{\R}{\mathbb{R}}
\newcommand{\C}{\mathbb{C}}
\newcommand{\F}{\mathbb{F}}
\newcommand{\A}{\mathbb{A}}
\newcommand{\PP}{\mathbb{P}}
\newcommand{\E}{\mathbb{E}}
\newcommand{\eps}{\varepsilon}
\newcommand{\dd}{\,\mathrm{d}}
\begin{document}
\section*{UnsolvedMath ID: 800265}
\textbf{Source problem code:} catalogue entry
\subsection{Definitions and mathematical statement}
Under an explicit conventional interpretation, a generalized complex Hénon map is the polynomial automorphism
\[
H_{p,a}\colon\C^2\longrightarrow\C^2,\qquad H_{p,a}(z,w)=(p(z)-aw,z),
\]
where \(p\in\C[z]\), \(\deg p\geq2\), and \(a\in\C\setminus\{0\}\). Write \(H^0=\operatorname{id}_{\C^2}\) and let \(H^n\) denote the \(n\)-fold iterate for integers \(n\geq1\).

For this formulation, use the forward Fatou set
\[
F^+(H)=\{x\in\C^2:\text{there is an open neighborhood }V\ni x\text{ on which }\{H^n\}_{n\geq0}\text{ is normal}\}.
\]
Here normality is taken after the standard inclusion \(\C^2\hookrightarrow\PP^2(\C)\): every sequence of these maps has a subsequence converging uniformly on compact subsets of \(V\) to a holomorphic map into \(\PP^2(\C)\). A forward Fatou component is a nonempty connected component \(U\) of \(F^+(H)\). It is wandering if its forward images are pairwise distinct Fatou components, that is,
\[
H^m(U)\neq H^n(U)\qquad\text{for all integers }0\leq m<n.
\]

With these conventions, the question is whether there exist \(p\in\C[z]\) of degree at least \(2\), \(a\in\C\setminus\{0\}\), and a forward Fatou component \(U\) of \(H_{p,a}\) such that \(U\) is wandering.

The supplied source asks only whether wandering domains exist for complex Hénon maps. It does not specify whether this means quadratic Hénon maps, the generalized maps displayed above, or finite compositions of generalized Hénon maps. It also does not supply a definition of wandering domain or the intended normality convention. Thus the displayed statement is a precise formulation under the stated conventions; the intended map class and dynamical conventions require human confirmation.
\subsection{Short English statement}
Does some complex Hénon map have a connected forward Fatou component whose forward images never repeat? The precise formulation above uses generalized Hénon maps and projective normality; the supplied record does not specify these conventions.
\subsection{Sources}
\begin{itemize}
\item Local UnsolvedMath registry, permanent ID \(800265\), source problem code \texttt{catalogue entry}, title: \emph{Several Complex Variables Problem List --- Problem 64}. The registry attributes the problem to Buzzard, Fornaess, and Krantz, with proposed year 1995. Its \texttt{open} label is supplied metadata, not an independently verified mathematical or literature assessment.
\item Buzzard, Fornaess, and Krantz, \emph{The Several Complex Variables Problem List} (1995), Problem 64. Supplied source URL: \url{https://arxiv.org/abs/math/9509202}. The supplied background reports access on 2026-10-16 and extraction by \texttt{arxiv\_latex}; these are source-provenance claims supplied with the record and were not independently verified here.
\end{itemize}
\end{document}
